%% toc: P15 158 158.4,158.5
% Proofs for Hahn partitions up to signed-null sets and Stieltjes representations (Durrett, Appendix A.4, p. 468).
\section{Hahn partitions and Stieltjes representations: proofs}
\label{sec:annex-p15-w02b-p15-005-hahn-stieltjes}

Use the definitions of
Chapter~\ref{sec:p15-w02b-p15-005-hahn-stieltjes} and
\cref{sec:p15-radon-nikodym}. A set is signed-null exactly when it is
both positive and negative, since a value is zero exactly when it is
both nonnegative and nonpositive. Consequently every measurable subset
of a signed-null set is signed-null, and by
\cref{thm:p15-signed-set-closure}, applied once to positivity and once to
negativity, a countable union of signed-null sets is signed-null.

\begin{annexitem}{annex:p15-w02b-p15-005-hahn-partitions}{thm:p15-w02b-p15-005-hahn-partitions}{Signed-null modifications of a Hahn partition}
\begin{proof}
(1) Write \(B'=A'^c\). The symmetric difference splits as
\[
  A\mathbin{\triangle}A'=(A'\setminus A)\cup(A\setminus A')
                        =(A'\cap B)\cup(A\cap B').
\]
Suppose first that \((A',B')\) is a Hahn partition. The set
\(A'\cap B\) is a measurable subset of the positive set \(A'\) and of
the negative set \(B\), so it is positive and negative by
\cref{thm:p15-signed-set-closure}, hence signed-null. The same argument
applies to \(A\cap B'\). Their union \(A\mathbin{\triangle}A'\) is
signed-null by the preliminary remark.

Conversely, suppose \(N=A\mathbin{\triangle}A'\) is signed-null. Then
\[
  A'=(A'\cap A)\cup(A'\cap B).
\]
The first piece is a measurable subset of the positive set \(A\); the
second is contained in \(A'\setminus A\subseteq N\), so it is
signed-null and in particular positive. A union of two positive sets is
positive, so \(A'\) is positive. Likewise
\[
  B'=(B'\cap B)\cup(B'\cap A),
\]
where the first piece lies in the negative set \(B\) and the second lies
in \(A\setminus A'\subseteq N\) and is therefore negative. Hence \(B'\)
is negative, and \((A',B')\) is a Hahn partition.

(2) By \cref{thm:p15-hahn-decomposition}, \(\alpha\) has at least one
Hahn partition \((A,B)\). If the only signed-null measurable set is
\(\varnothing\), then any Hahn positive set \(A'\) satisfies
\(A\mathbin{\triangle}A'=\varnothing\) by part~(1), so \(A'=A\) and the
partition is unique. If instead \(N\neq\varnothing\) is measurable and
signed-null, put \(A'=A\mathbin{\triangle}N\). It is measurable and
\(A\mathbin{\triangle}A'=N\), because symmetric difference is
associative and \(A\mathbin{\triangle}A=\varnothing\). By part~(1),
\((A',A'^c)\) is a Hahn partition, and \(A'\neq A\) since their
symmetric difference is the nonempty set \(N\).

(3) Let \(E\) be measurable and \(N=A\mathbin{\triangle}A'\), which is
signed-null by part~(1). The disjoint decompositions
\[
  E\cap A'=(E\cap A\cap A')\cup(E\cap(A'\setminus A)),\qquad
  E\cap A=(E\cap A\cap A')\cup(E\cap(A\setminus A'))
\]
have second pieces contained in \(N\), so those pieces have value zero.
Finite additivity gives
\(\alpha(E\cap A')=\alpha(E\cap A\cap A')=\alpha(E\cap A)\).
Since \(A'^c\mathbin{\triangle}B=A'^c\mathbin{\triangle}A^c
=A\mathbin{\triangle}A'\), the same argument with complements gives
\(\alpha(E\cap A'^c)=\alpha(E\cap B)\). By
\cref{thm:p15-jordan-decomposition}, both partitions therefore give the
same Jordan components.
\end{proof}
\end{annexitem}

\begin{annexitem}{annex:p15-w02b-p15-005-density-hahn}{thm:p15-w02b-p15-005-density-hahn}{The density criterion for Hahn positive sets}
\begin{proof}
The sets \(\{f<0\}\), \(\{f=0\}\) and \(\{f>0\}\) are measurable
because \(f\) is measurable. By \cref{thm:p15-signed-measure-examples},
a measurable set \(D\) is positive exactly when
\(\mu(D\cap\{f<0\})=0\) and negative exactly when
\(\mu(D\cap\{f>0\})=0\).

(1) The set \(N\) is signed-null exactly when it is positive and
negative, that is, exactly when
\(\mu(N\cap\{f<0\})=0\) and \(\mu(N\cap\{f>0\})=0\). Since
\(N\cap\{f\neq0\}\) is the disjoint union of these two sets, additivity
of \(\mu\) shows that both vanish exactly when
\(\mu(N\cap\{f\neq0\})=0\); one direction uses that each of the two
sets is contained in \(N\cap\{f\neq0\}\).

(2) The pair \((A,A^c)\) is a Hahn partition exactly when \(A\) is
positive and \(A^c\) is negative. By the criterion above, this means
\(\mu(A\cap\{f<0\})=0\) and \(\mu(A^c\cap\{f>0\})=0\). Since
\(\{f\geq0\}^c=\{f<0\}\), the first condition says
\(A\subseteq\{f\geq0\}\) \(\mu\)-almost everywhere in the sense of
\cref{def:p15-w02b-p15-005-hahn-partition}, and the second says
\(\{f>0\}\subseteq A\) \(\mu\)-almost everywhere. For \(A=\{f>0\}\),
the sets \(A\cap\{f<0\}\) and \(A^c\cap\{f>0\}\) are empty. For
\(A=\{f\geq0\}\), the set \(A\cap\{f<0\}\) is empty and
\(A^c\cap\{f>0\}=\{f<0\}\cap\{f>0\}\) is empty. Both are therefore Hahn
positive sets.

(3) Here \(f\) is continuous, hence Borel, and
\(\int f^-\,d\lambda=\int_0^{1/2}(1-2x)\,dx=1/4<\infty\), so part~(2)
applies, with \(\{f>0\}=(1/2,1]\) and \(\{f<0\}=[0,1/2)\). Let
\(Z\subseteq[0,1]\) be Borel with \(\lambda(Z)=0\).
For \(A=(1/2,1]\cup Z\), the set \(\{f>0\}\cap A^c\) is empty, and
\(A\cap\{f<0\}\subseteq Z\) has measure zero.
For \(A=(1/2,1]\setminus Z\), the set \(\{f>0\}\cap A^c\) equals
\((1/2,1]\cap Z\), of measure zero, and \(A\cap\{f<0\}\) is empty.
By part~(2) both are Hahn positive sets.

The singleton \(\{1/2\}\) is Borel with Lebesgue measure zero
(\cref{thm:p15-lebesgue-unit-interval}), giving
\((1/2,1]\cup\{1/2\}=[1/2,1]\). The Cantor set \(C\) is a Borel subset
of \([0,1]\) with \(\lambda(C)=0\) by
\cref{thm:p15-cantor-singular-law}. It contains \(0\): for every
\(n\), the all-zero word \(b=(0,\ldots,0)\) has \(c_b=0\), so
\(0\in I_b\subseteq C_n\), and hence \(0\in\bigcap_nC_n=C\).
The three sets \((1/2,1]\), \([1/2,1]\) and \((1/2,1]\cup C\) are
distinct: the point \(1/2\) lies in the second and not the first, and
the point \(0\) lies in the third and in neither of the others. Finally
\(f(0)=-1\).
\end{proof}
\end{annexitem}

\begin{annexitem}{annex:p15-w02b-p15-005-stieltjes-excess}{thm:p15-w02b-p15-005-stieltjes-excess}{The excess measure and the minimality of the Jordan pair}
\begin{proof}
(1) Each of \(F,G\) is a Stieltjes measure function in the sense of
\cref{def:p15-stieltjes-intervals}, so
\cref{thm:p15-stieltjes-construction} gives Borel measures
\(\mu_F,\mu_G\) with \(\mu_F(I(a,b))=F(b)-F(a)\), where
\(F(-\infty)=\inf F=0\) because \(F\) is nondecreasing with limit zero
at \(-\infty\). Hence \(\mu_F((-\infty,x])=F(x)\) for real \(x\), and
\(\mu_F(\mathbb R)=F(\infty)<\infty\) because \(F\) is bounded. The
same holds for \(G\). By \cref{thm:p15-signed-measure-examples}, the
difference \(\mu_F-\mu_G\) is a signed measure; both terms are finite,
so it is a finite signed Borel measure, and its cumulative function is
\(F-G\). The finite signed measure \(\alpha\) has the same cumulative
function by hypothesis, so the uniqueness assertion of
\cref{thm:p15-signed-cdf-variation} gives \(\alpha=\mu_F-\mu_G\).

(2) The set functions \(E\mapsto\mu_F(E\cap B)\) and
\(E\mapsto\mu_G(E\cap A)\) are finite Borel measures: intersecting a
disjoint family with a fixed Borel set keeps it disjoint, and countable
additivity of \(\mu_F\) and \(\mu_G\) passes to the restrictions. Their
sum \(\rho\) is a finite Borel measure. For Borel \(E\), part~(1) and
\(E=(E\cap A)\cup(E\cap B)\) give
\[
 \begin{aligned}
 \mu_F(E)-\alpha^+(E)
   &=\mu_F(E\cap A)+\mu_F(E\cap B)
      -\bigl(\mu_F(E\cap A)-\mu_G(E\cap A)\bigr)\\
   &=\rho(E),\\
 \mu_G(E)-\alpha^-(E)
   &=\mu_G(E\cap A)+\mu_G(E\cap B)
      +\bigl(\mu_F(E\cap B)-\mu_G(E\cap B)\bigr)\\
   &=\rho(E),
 \end{aligned}
\]
using \(\alpha^+(E)=\alpha(E\cap A)\) and \(\alpha^-(E)=-\alpha(E\cap B)\)
from \cref{thm:p15-jordan-decomposition}. All terms are finite, so the
subtractions are legitimate, and the two displayed identities are
\(\mu_F=\alpha^++\rho\) and \(\mu_G=\alpha^-+\rho\). If
\(\mu_F=\alpha^++\rho'\) for some measure \(\rho'\), then
\(\rho'(E)=\mu_F(E)-\alpha^+(E)=\rho(E)\) for every \(E\), since
\(\alpha^+(E)\) is finite; the same holds for the second identity. As
\(\alpha^+\) and \(\alpha^-\) do not depend on the Hahn partition, neither
does \(\rho\). Nonnegativity of \(\rho\) gives the two inequalities.
Evaluating the identities at \(\mathbb R\) and adding gives
\[
 F(\infty)+G(\infty)=\mu_F(\mathbb R)+\mu_G(\mathbb R)
  =\alpha^+(\mathbb R)+\alpha^-(\mathbb R)+2\rho(\mathbb R)
  =\lvert\alpha\rvert(\mathbb R)+2\rho(\mathbb R).
\]

(3) By part~(2), \(\mu_F=\alpha^+\) holds exactly when \(\rho=0\), and
then also \(\mu_G=\alpha^-\); conversely \(\mu_G=\alpha^-\) and
\(\mu_F=\alpha^+\) force \(\rho=0\). If \(\rho=0\), then
\(\mu_F=\alpha^+\) and \(\mu_G=\alpha^-\) are mutually singular by
\cref{thm:p15-jordan-decomposition}. If \(\mu_F\perp\mu_G\), then
\(\mu_F,\mu_G\) are mutually singular measures with
\(\alpha=\mu_F-\mu_G\) and \(\mu_G(\mathbb R)<\infty\), so the uniqueness
in \cref{thm:p15-jordan-decomposition} gives \(\mu_F=\alpha^+\) and
\(\mu_G=\alpha^-\). Finally, \(\rho(\mathbb R)\) is finite and
nonnegative, so by the mass identity of part~(2),
\(F(\infty)+G(\infty)=\lvert\alpha\rvert(\mathbb R)\) holds exactly
when \(\rho(\mathbb R)=0\), that is, exactly when \(\rho=0\).

(4) For a finite Borel measure \(\tau\), the function
\(x\mapsto\tau((-\infty,x])\) is bounded by \(\tau(\mathbb R)\),
nonnegative and nondecreasing; it is right continuous by continuity
from above along \((-\infty,x+1/n]\), which uses finiteness of \(\tau\);
and it tends to zero at \(-\infty\) by continuity from above along
\((-\infty,-n]\), whose intersection is empty. Apply this to the finite
measures \(\alpha^++\rho\) and \(\alpha^-+\rho\). The difference of the
two functions at \(x\) is
\(\alpha^+((-\infty,x])-\alpha^-((-\infty,x])=\alpha((-\infty,x])\),
because the finite value \(\rho((-\infty,x])\) cancels.

(5) If \(F=G=\mathbf1_{[0,\infty)}\), both functions have every
required property and \(F-G=0\). The zero signed measure has cumulative
function zero, so part~(1) gives \(\alpha=0\). For \(\alpha=0\) the
pair \((\mathbb R,\varnothing)\) is a Hahn partition, so
\(\alpha^+=\alpha^-=0\). Here \(F(-\infty)=0\) and \(F(\infty)=1\).
For all allowed endpoints \(-\infty\leq a<b\leq\infty\),
\(F(b)-F(a)\) equals \(1\) when \(a<0\leq b\) and \(0\) otherwise,
which is \(\delta_0(I(a,b))\). By the uniqueness in
\cref{thm:p15-stieltjes-construction}, \(\mu_F=\delta_0\), and likewise
\(\mu_G=\delta_0\). Part~(2) then gives
\(\rho=\mu_F-\alpha^+=\delta_0\).

For \(F(x)=2\min(\max(x,0),1)\) and \(G(x)=\min(\max(x,0),1)\), both
functions are continuous, nonnegative, nondecreasing, bounded and zero
on \((-\infty,0]\), with \(G(\infty)=1\) and \(F(\infty)=2\). For all
allowed endpoints \(-\infty\leq a<b\leq\infty\), the difference
\(G(b)-G(a)\) is the length of the interval \(I(a,b)\cap[0,1]\), which
is \(\lambda(I(a,b)\cap[0,1])\) by \cref{thm:p15-lebesgue-unit-interval},
and \(F(b)-F(a)\) is twice that. Uniqueness in
\cref{thm:p15-stieltjes-construction} gives
\(\mu_G=\lambda(\,\cdot\,\cap[0,1])\) and \(\mu_F=2\mu_G\), so by
part~(1), \(\alpha=\mu_F-\mu_G=\lambda(\,\cdot\,\cap[0,1])\). This
\(\alpha\) is a measure, so \((\mathbb R,\varnothing)\) is a Hahn
partition, \(\alpha^+=\alpha\) and \(\alpha^-=0\), and part~(2) gives
\(\rho=\mu_F-\alpha^+=\alpha\).

(6) Let \(F\) have the required properties, not be identically zero,
and put \(G=F\). Then \(F-G=0\) is the cumulative function of the zero
signed measure, so part~(1) gives \(\alpha=\mu_F-\mu_G=0\). As in
part~(5), \((\mathbb R,\varnothing)\) is a Hahn partition for \(\alpha=0\),
so \(\alpha^+=\alpha^-=0\), and part~(2) gives
\(\rho=\mu_F-\alpha^+=\mu_F\) and \(\rho=\mu_G-\alpha^-=\mu_G\).
Choose \(x\) with \(F(x)\neq0\); by part~(1),
\(\mu_F((-\infty,x])=F(x)>0\), because \(F\) is nonnegative. Thus
\(\mu_F=\mu_G=\rho\) is a nonzero finite measure, while both Jordan
components are zero, so \(\mu_F\neq\alpha^+\) and \(\mu_G\neq\alpha^-\).
\end{proof}
\end{annexitem}
